\documentclass[11pt]{article}
\usepackage[a4paper,margin=24mm]{geometry}
\usepackage{amsmath,amssymb,hyperref}
\newcommand{\dd}{\mathrm{d}}
\title{Four-loop two-point graphs with one weighted bubble edge}
\author{Phi3Lab: isolated continuation}
\date{25 October 2026}
\begin{document}\maketitle
The scope is L4E2G4, L4E2G13, L4E2G25 and L4E2G47.
All formulas use Euclidean dimension six, $m^2=p^2=1$, and
$\dd^6 k/\pi^3$ per loop. Couplings and graph symmetry weights are excluded,
as in the frozen original numerical programs. The accepted lower-loop
L3E2G5 is not rerun.

\section*{Exact forest pairing}
The original 36 compatible forests pair into 18 states after combining the
one-loop two-point bubble with its degree-two local subtraction. All
overlapping vertex subtractions remain in the paired sum. The two original
bubble-adjacent propagators are retained, giving
\[
 W(t)=\frac{B(t)}{(1+t)^2},\qquad
 B(t)=\frac{(t+4)^{3/2}}{3\sqrt t}
       \operatorname{arsinh}\frac{\sqrt t}{2}
       -\frac{4t}{9}-\frac43 .
\]
Here $t=k^2$ is a retained host-loop invariant. In particular, a parent
$T^0$ on physical external momenta does not set $k$ to zero.
Additional edges can acquire the same route after a vertex projection;
they remain separate edges. The independent spectral identity is
\[
 B(t)=t^2\int_0^{1/4}\frac{(1-4\sigma)^{3/2}}{6(1+\sigma t)}\,\dd\sigma .
\]
Both $B(0)$ and $B'(0)$ vanish. Derivatives with respect to $t=p^2$ remove
the constant and $p^2$ terms, not a $p^4$ term.

\section*{Analytic common scale}
Each paired forest has three effective loops and eight effective edges.
For its routed quadratic form let $U_f$ be the loop determinant,
$F_f$ the augmented determinant, and $Q_f=F_f/U_f$.
Positive Cauchy--Binet minors provide an independent construction.
With $M=\sum_{e=1}^8\alpha_e$ and $d$ the weighted-edge parameter, the
complete projective density is
\[
 \mathcal K(\alpha)=
 \sum_{f\in\mathcal F_{\rm paired}}(-1)^{|f|}
 \frac{J(M,d,Q_f)}{U_f^3}.
\]
The common Schwinger scale has power $\lambda^{-2}$. The overall
degree-two subtraction is retained in the initial-value problem
\[
 J(0)=J'(0)=0,\qquad
 J''(z)=\frac{d\,B((M+z-d)/d)}{(M+z)^2}.
\]
Its full independent remainder integral is
\[
 J(M,d,z)=d z^2\int_0^1
 \frac{(1-u)B((M+uz-d)/d)}{(M+uz)^2}\,\dd u .
\]
The whole signed forest sum is evaluated at each shared point.
No claim is made that individual forest terms have separately convergent
parameter integrals.

Let $A'(t)=B(t)/(1+t)^2$, $A(0)=0$, using the independently verified
logarithm/dilogarithm primitive. Put $h=\operatorname{arsinh}(\sqrt t/2)$ and
\[
 \begin{split}
 C(t)={}&\frac{h\sqrt{t(t+4)}(t-2)}{3(t+1)}
       +\frac43h^2-\frac{11}{18}t\\
       &+\frac1{18}\log(1+t)+\frac{8t}{9(t+1)} .
 \end{split}
\]
Then $C'(t)=tB(t)/(1+t)^2$ and $C(0)=0$. The closed scale primitive is
\[
 t_0=\frac{M-d}{d},\qquad t_1=t_0+\frac zd,\qquad
 J=d\left\{t_1[A(t_1)-A(t_0)]-[C(t_1)-C(t_0)]\right\}.
\]
These derivative and boundary identities are checked separately from
numerical comparisons.

\section*{Eight exact sectors}
The degrees of $J,U_f,F_f$ are $1,3,4$, so $\mathcal K$ is homogeneous
of degree $-8$. Partition the positive parameter domain according to its
largest coordinate. In sector $i$, set $\alpha_i=1$ and let the remaining
seven coordinates $x_j$ range over $(0,1)$. Normalizing by
$S=1+\sum_j x_j$ gives Jacobian $S^{-8}$, exactly cancelling
$\mathcal K(\alpha/S)=S^8\mathcal K(\alpha)$. Therefore
\[
 \int_{\Delta_7}\mathcal K(\alpha)\,\dd^7\alpha
 =\sum_{i=1}^8\int_{(0,1)^7}
   \mathcal K(x_1,\ldots,1,\ldots,x_7)\,\dd^7x .
\]
The eight sectors are disjoint up to null boundaries. No edge-permutation
symmetry is assumed and there is no extra factor $7!$.
The simplex volume is $1/5040$.
The predefined maps are $x_j=u_j$ and $x_j=u_j^2$; their Jacobians are
$1$ and $\prod_j2u_j$, respectively.

\section*{Evidence and limits}
Graph proof02 checks all 36 original forest identities. Measure proof01
independently checks 8188 nonempty edge subsets and all local Taylor orders.
Scalar proof03 compares six full remainder integrals at 320 digits,
requiring 200 relative digits, and adds twelve hypergeometric bubble checks.
Native values are Arb intervals; the independent reference constructs
original Gaussian matrices and uses a Schur solve.
Actual process exits and source hashes remain the authority for completion.
No graph is accepted from a symbolic proof or compilation alone.

Convergence uses the massive Euclidean BPH theorem together with the
exhaustive local checks; that theorem is not reproved by these scripts.
Empirical sampling acceptance is not a rigorous quadrature error
certificate. For the BPH framework see H.~van Hees,
\emph{Introduction to Relativistic Quantum Field Theory},
Sections 5.7--5.9,
\url{https://itp.uni-frankfurt.de/~hees/publ/lect.pdf}.
\end{document}